\subsection{Controller 5-3: Projection-Based Adaptive (Finite Time)}
\label{sec: controller5-3}

Controller 5-1 needs a bound on $\|Y_i\tilde{\Theta}_i\|$, while Controller 5-2 avoids any
a~priori gain specification at the cost of losing finite-time convergence. Controller 5-3
recovers finite-time sliding by projecting the parameter estimate onto a known compact
set, which converts the required bound into a bound on the true parameter vector
$\|\Theta_i\|$ --- a far more natural assumption for physical systems
\cite{ioannouSunRobustAdaptiveControl1996}.

\subsubsection{Assumption and Projection}

We replace the gain condition \eqref{eq5:rho} of Controller 5-1 with the following
two-part assumption. The first part is a standard physical bound; the second part
follows from the continuity of $Y_i$ (Assumption~A3) and the boundedness of its
arguments on the physically realisable domain of the EL system.

\begin{assumption}
	\hfill
	\begin{enumerate}
		\item[(i)] The true parameter vector satisfies $\|\Theta_i\| \le \Theta_{i,\max}$
			for some known $\Theta_{i,\max} > 0$.
		\item[(ii)] The regressor satisfies $\|Y_i(q_i,\dot{q}_i,\dot{\zeta}_i,\zeta_i)\| \le k_Y$
			for a constant $k_Y > 0$, uniformly in $t$.
	\end{enumerate}
\end{assumption}

\noindent Assumption~(ii) is justified by the following facts.  Assumption~A3 gives
$Y_i$ as a continuous function of its arguments.  The target trajectory is bounded
(Assumption~A4).  The state-dependent signals $\zeta_i$ and $\dot{\zeta}_i$ defined in
\eqref{eq5-3:zeta} are continuous functions of $q_{i0}$, $\dot{q}_{i0}$, $\phi_i$ and the
target states.  Physically, the generalised coordinates $q_i$ of an EL system always
lie in a bounded set (e.g.~robot joint limits), and the velocities $\dot{q}_i$ are
bounded on any finite time interval by the system's passivity.  Consequently, there
exists a finite $k_Y$ that bounds $\|Y_i\|$ on the compact set in which the system
evolves before collision; this bound is independent of the collision-avoidance
analysis and is solely a property of the physical plant and the target.

Define the projection operator
$\mathrm{Proj}_{\hat{\Theta}_i} : \mathbb{R}^p \to \mathbb{R}^p$ that keeps $\hat{\Theta}_i$
inside the ball $\mathcal{B}_i = \{\Theta \in \mathbb{R}^p : \|\Theta\| \le \Theta_{i,\max}\}$:
\begin{equation}\label{eq5-3:proj}
	\mathrm{Proj}_{\hat{\Theta}_i}(v) =
	\begin{cases}
		v, & \|\hat{\Theta}_i\| < \Theta_{i,\max}
		\;\text{or}\; \hat{\Theta}_i^T v \le 0,\\[2pt]
		v - \dfrac{\hat{\Theta}_i(\hat{\Theta}_i^T v)}{\Theta_{i,\max}^2},
		& \|\hat{\Theta}_i\| = \Theta_{i,\max}
		\;\text{and}\; \hat{\Theta}_i^T v > 0.
	\end{cases}
\end{equation}
The key property of this projection (see \cite{ioannouSunRobustAdaptiveControl1996}) is
\begin{equation}\label{eq5-3:proj_prop}
	\tilde{\Theta}_i^T \Lambda\,\mathrm{Proj}_{\hat{\Theta}_i}(-\Lambda^{-1}Y_i^T s_i)
	\le -\tilde{\Theta}_i^T Y_i^T s_i,
\end{equation}
which holds for all $\hat{\Theta}_i \in \mathcal{B}_i$ whenever $\Theta_i \in \mathcal{B}_i$.

\subsubsection{Integral APF, Sliding Variable and Control Law}

The integral APF and sliding variable are identical to those of Controller 5-1:
\begin{equation}\label{eq5-3:zeta}
	\zeta_i = \dot{q}_0 - k_\alpha q_{i0}
	+ \int_0^t \bigl(-k_\alpha q_{i0} + \phi_i(\tau)\bigr)\,d\tau,\qquad
	s_i = \dot{q}_i - \zeta_i.
\end{equation}
The control input is
\begin{equation}\label{eq5-3:control_input}
	\tau_i = -k s_i - \rho_i\,\mathrm{sgn}(s_i)
	+ Y_i(q_i,\dot{q}_i,\dot{\zeta}_i,\zeta_i) \hat{\Theta}_i,\qquad k>0,
\end{equation}
with the robust gain chosen as
\begin{equation}\label{eq5-3:rho}
	\rho_i = k_Y \cdot 2\Theta_{i,\max} + \eta_i,\qquad \eta_i > 0.
\end{equation}
The regression equation is the same as in \eqref{eq5:regression}, and the adaptation law
is the projected version
\begin{equation}\label{eq5-3:adaptation}
	\dot{\hat{\Theta}}_i = \mathrm{Proj}_{\hat{\Theta}_i}\bigl(-\Lambda^{-1}Y_i^T s_i\bigr).
\end{equation}

\subsubsection{Closed-Loop Dynamics and Finite-Time Sliding}

Substituting \eqref{eq5-3:control_input} into \eqref{eq hetero agent} gives
\begin{equation}\label{eq5-3:s_dynamics}
	M_i(q_i)\dot{s}_i + C_i(q_i,\dot{q}_i)s_i
	= -k s_i - \rho_i\,\mathrm{sgn}(s_i)
	+ Y_i(q_i,\dot{q}_i,\dot{\zeta}_i,\zeta_i)\tilde{\Theta}_i.
\end{equation}
Consider $V_{1i} = \frac12 s_i^T M_i s_i + \frac12 \tilde{\Theta}_i^T \Lambda \tilde{\Theta}_i$.
Differentiating and using \eqref{eq5-3:proj_prop},
\begin{align}
	\dot{V}_{1i}
	&= -k\|s_i\|^2 - \rho_i\|s_i\|_1
		+ s_i^T Y_i\tilde{\Theta}_i
		+ \tilde{\Theta}_i^T \Lambda\,\dot{\hat{\Theta}}_i \notag\\
	&\le -k\|s_i\|^2 - \rho_i\|s_i\|_1
		+ s_i^T Y_i\tilde{\Theta}_i
		- \tilde{\Theta}_i^T Y_i^T s_i \notag\\
	&= -k\|s_i\|^2 - \rho_i\|s_i\|_1 \le 0. \label{eq5-3:V1_dot}
\end{align}
For the sliding part $V_{si} = \frac12 s_i^T M_i s_i$ alone,
\begin{equation}\label{eq5-3:Vs_dot}
	\dot{V}_{si} = -k\|s_i\|^2 - \rho_i\|s_i\|_1 + s_i^T Y_i\tilde{\Theta}_i
	\le -k\|s_i\|^2 - \eta_i\|s_i\|_1
	\le -\eta_i\|s_i\|,
\end{equation}
where we used $\|s_i\|_1 \ge \|s_i\|_2$ and
$\|s_i^T Y_i\tilde{\Theta}_i\| \le k_Y \cdot 2\Theta_{i,\max}\,\|s_i\| \le \rho_i\|s_i\|$.
With $V_{si} \le \tfrac{k_{\overline m}}{2}\|s_i\|^2$, this yields
$\dot{V}_{si} \le -c_i\sqrt{V_{si}}$ with $c_i = \eta_i\sqrt{2/k_{\overline m}}$, so
$s_i \equiv 0$ for all $t \ge T^* = \max_i T_i$ with finite $T_i$. Hence
$\dot{s}_i \equiv 0$ for $t \ge T^*$, exactly as in Controller 5-1.

\subsubsection{Theoretical Guarantee}

\begin{theorem}
	Let Assumptions A1--A3 and the projection assumption above hold. The HELS
	\eqref{eq hetero agent} governed by \eqref{eq5-3:control_input} and
	\eqref{eq5-3:adaptation} achieves the fencing objectives P1--P3 asymptotically,
	without the PE condition and without differentiating $\phi_i$.
\end{theorem}

\noindent\textbf{Proof}

The proof is organized into eight steps. Steps 1--2 establish the finite-time convergence of the
sliding variable $s_i$ using the projection operator; Steps 3--4 analyze the boundedness and
collision avoidance during the transient interval $[0,T^*]$; Steps 5--8 exploit the post-sliding
clean dynamics to prove P3, P1 and P2 asymptotically.

\medskip\noindent\textbf{Step 1 (Finite-time convergence of $s_i$).}

Differentiating $V_{1i} = \frac12 s_i^T M_i s_i + \frac12 \tilde{\Theta}_i^T \Lambda \tilde{\Theta}_i$
and using \eqref{eq5-3:s_dynamics}, \eqref{eq5-3:adaptation} and the skew-symmetry property A2,
\begin{align}
	\dot{V}_{1i}
	&= s_i^T M_i \dot{s}_i + \tfrac12 s_i^T \dot{M}_i s_i
		+ \tilde{\Theta}_i^T \Lambda \dot{\tilde{\Theta}}_i \notag\\
	&= s_i^T\bigl(-k s_i - C_i s_i - \rho_i\,\mathrm{sgn}(s_i) + Y_i\tilde{\Theta}_i\bigr)
		+ \tfrac12 s_i^T \dot{M}_i s_i
		- \tilde{\Theta}_i^T \Lambda \dot{\hat{\Theta}}_i \notag\\
	&= -k\|s_i\|^2 - \rho_i\|s_i\|_1 + s_i^T Y_i\tilde{\Theta}_i
		+ \tfrac12 s_i^T(\dot{M}_i - 2C_i)s_i
		- \tilde{\Theta}_i^T \Lambda \dot{\hat{\Theta}}_i. \label{eq5-3:V1_dot_raw}
\end{align}
The skew-symmetry property A2 gives $\tfrac12 s_i^T(\dot{M}_i - 2C_i)s_i = 0$.
Using the projection adaptation \eqref{eq5-3:adaptation} and the projection property
\eqref{eq5-3:proj_prop},
\begin{equation}
	-\tilde{\Theta}_i^T \Lambda \dot{\hat{\Theta}}_i
	= -\tilde{\Theta}_i^T \Lambda\,\mathrm{Proj}_{\hat{\Theta}_i}\bigl(-\Lambda^{-1}Y_i^T s_i\bigr)
	\le \tilde{\Theta}_i^T Y_i^T s_i = s_i^T Y_i\tilde{\Theta}_i. \label{eq5-3:proj_ineq}
\end{equation}
Substituting \eqref{eq5-3:proj_ineq} into \eqref{eq5-3:V1_dot_raw} yields
\begin{equation}
	\dot{V}_{1i} \le -k\|s_i\|^2 - \rho_i\|s_i\|_1 + 2 s_i^T Y_i\tilde{\Theta}_i. \label{eq5-3:V1_dot_proj}
\end{equation}
\noindent Unlike Controller 5-1, the cross term $s_i^T Y_i\tilde{\Theta}_i$ does not cancel
exactly under projection; the projection property only provides an inequality.
Nevertheless, the robust gain $\rho_i$ is chosen large enough to dominate this term.

Now consider the sliding part $V_{si} = \frac12 s_i^T M_i s_i$ alone.
Differentiating and using the skew-symmetry,
\begin{align}
	\dot{V}_{si}
	&= s_i^T M_i \dot{s}_i + \tfrac12 s_i^T \dot{M}_i s_i \notag\\
	&= s_i^T\bigl(-k s_i - C_i s_i - \rho_i\,\mathrm{sgn}(s_i) + Y_i\tilde{\Theta}_i\bigr)
		+ \tfrac12 s_i^T \dot{M}_i s_i \notag\\
	&= -k\|s_i\|^2 - \rho_i\|s_i\|_1 + s_i^T Y_i\tilde{\Theta}_i. \label{eq5-3:Vs_dot}
\end{align}
The projection property is not needed here because $\dot{V}_{si}$ does not contain the
adaptation term $\tilde{\Theta}_i^T \Lambda \dot{\hat{\Theta}}_i$.

By Assumption~(ii), $\|\Theta_i\| \le \Theta_{i,\max}$ and $\|Y_i\| \le k_Y$ hold for all $t$
(see the justification below the assumption).  The projection \eqref{eq5-3:proj} ensures
$\|\hat{\Theta}_i\| \le \Theta_{i,\max}$ for all $t$, hence
$\|\tilde{\Theta}_i\| \le \|\hat{\Theta}_i\| + \|\Theta_i\| \le 2\Theta_{i,\max}$.
Therefore,
\begin{equation}
	\|s_i^T Y_i\tilde{\Theta}_i\|
	\le \|s_i\| \, \|Y_i\| \, \|\tilde{\Theta}_i\|
	\le \|s_i\| \cdot k_Y \cdot 2\Theta_{i,\max}. \label{eq5-3:Ytb_bound}
\end{equation}
Since $\rho_i = k_Y \cdot 2\Theta_{i,\max} + \eta_i$ with $\eta_i > 0$, we have
\begin{equation}
	-\rho_i\|s_i\|_1 + s_i^T Y_i\tilde{\Theta}_i
	\le -\rho_i\|s_i\| + \|s_i\| \, k_Y \cdot 2\Theta_{i,\max}
	\le -\eta_i\|s_i\|, \label{eq5-3:Vs_bound}
\end{equation}
where we used $\|s_i\|_1 \ge \|s_i\|_2 = \|s_i\|$.
Substituting \eqref{eq5-3:Vs_bound} into \eqref{eq5-3:Vs_dot} gives
\begin{equation}
	\dot{V}_{si} \le -k\|s_i\|^2 - \eta_i\|s_i\|
	\le -\eta_i\|s_i\|. \label{eq5-3:Vs_dot_bound}
\end{equation}

From A1, $M_i(q_i) \le k_{\overline m} I_p$, so $V_{si} \le \frac{k_{\overline m}}{2}\|s_i\|^2$,
which implies $\|s_i\| \ge \sqrt{2V_{si}/k_{\overline m}}$.  Combining with \eqref{eq5-3:Vs_dot_bound},
\begin{equation}
	\dot{V}_{si} \le -\eta_i \sqrt{\frac{2}{k_{\overline m}}} \sqrt{V_{si}}
	= -c_i \sqrt{V_{si}}, \qquad
	c_i = \eta_i \sqrt{\frac{2}{k_{\overline m}}}. \label{eq5-3:finite_time_ode}
\end{equation}
The differential inequality $\dot{V}_{si} \le -c_i \sqrt{V_{si}}$ implies that $V_{si}$ reaches
zero in finite time (comparison lemma, see e.g.~\cite[Lemma~3.6]{khalilNonlinearSystems2002}):
\begin{equation}
	T_i \le \frac{2\sqrt{V_{si}(0)}}{c_i},
	\qquad
	s_i(t) \equiv 0,\; \dot{s}_i(t) \equiv 0,\; \forall\, t \ge T_i. \label{eq5-3:finite_time}
\end{equation}
Define $T^* = \max_{i\in\mathcal{N}} T_i$; then for all $t \ge T^*$, $s_i(t) \equiv 0$ and
$\dot{s}_i(t) \equiv 0$ for every $i \in \mathcal{N}$.  No PE condition is invoked anywhere
in this step.

\medskip\noindent\textbf{Step 2 (Boundedness of $\dot{V}_{1i}$ and parameter estimates).}

From \eqref{eq5-3:V1_dot_proj} and the bound $\|s_i^T Y_i\tilde{\Theta}_i\| \le \rho_i\|s_i\|$,
\begin{equation}
	\dot{V}_{1i} \le -k\|s_i\|^2 - \rho_i\|s_i\|_1 + 2\rho_i\|s_i\|
	\le -k\|s_i\|^2 + \rho_i\|s_i\|. \label{eq5-3:V1i_bound}
\end{equation}
While $s_i \neq 0$, the right-hand side may be positive when $\|s_i\|$ is small.
Nevertheless, $V_{1i}$ can grow at most linearly on any finite interval, and once $s_i \equiv 0$,
$\dot{V}_{1i} \le 0$.  Hence $V_{1i}(t)$ is bounded on $[0,T^*]$, implying $\tilde{\Theta}_i$ is
bounded for all $t \ge 0$.

\medskip\noindent\textbf{Step 3 (No finite escape time on $[0,T^*]$).}

On the interval $[0,T^*]$, the sliding dynamics \eqref{eq5-3:s_dynamics} give
\begin{equation}
	M_i \dot{s}_i = -k s_i - C_i s_i - \rho_i\,\mathrm{sgn}(s_i) + Y_i\tilde{\Theta}_i.
\end{equation}
Since $s_i$ converges to zero in finite time (hence bounded on $[0,T^*]$), $\tilde{\Theta}_i$ is
bounded (Step 2), and $Y_i$ is bounded by the projection assumption, the right-hand side is
uniformly bounded.  By A1, $M_i$ is invertible with $M_i^{-1} \le \frac{1}{k_{\underline m}} I_p$,
so $\dot{s}_i$ is uniformly bounded on $[0,T^*]$.  Consequently, $\dot{s}_i$ is bounded and
$s_i$ is Lipschitz, so no finite escape occurs.

\medskip\noindent\textbf{Step 4 (Collision avoidance during the transient $[0,T^*]$).}

From the definition of $s_i$ and $\dot{s}_i$,
\begin{equation}
	\ddot{q}_{i0} = -k_\alpha \dot{q}_{i0} - k_\alpha q_{i0} + \phi_i + \dot{s}_i. \label{eq5-3:qddot_transient}
\end{equation}
Consider the fencing Lyapunov function
\begin{equation}
	V_2 = \frac12\sum_{i\in\mathcal{N}}\sum_{j\in\mathcal{N}}
	\int_{\|q_{ij}\|}^{\mu}\alpha(s)\,ds
	+ \frac{k_\alpha}{2}\sum_{i\in\mathcal{N}}\|q_{i0}\|^2
	+ \frac12\sum_{i\in\mathcal{N}}\|\dot{q}_{i0}\|^2. \label{eq5-3:V2}
\end{equation}
Its time derivative, using $\sum_i\phi_i = 0$ and $\dot{A} = -\sum_i\phi_i^T\dot{q}_{i0}$, is
\begin{align}
	\dot{V}_2
	&= -\sum_i\phi_i^T\dot{q}_{i0}
		+ k_\alpha\sum_i q_{i0}^T\dot{q}_{i0}
		+ \sum_i \dot{q}_{i0}^T\ddot{q}_{i0} \notag\\
	&= -k_\alpha\sum_i\|\dot{q}_{i0}\|^2
		+ \sum_i \dot{q}_{i0}^T\dot{s}_i. \label{eq5-3:V2_dot_transient}
\end{align}
Let $\bar\sigma = \max_{i\in\mathcal{N}} \sup_{t\in[0,T^*]} \|\dot{s}_i(t)\|$, which is finite by
Step 3.  Then
\begin{equation}
	\dot{V}_2 \le \frac{k_\alpha}{2}\sum_i\|\dot{q}_{i0}\|^2
		+ \frac{N\bar\sigma^2}{2k_\alpha}. \label{eq5-3:V2_dot_transient_bound}
\end{equation}
Integrating over $[0,t]$ for any $t \le T^*$,
\begin{equation}
	V_2(t) \le V_2(0) + \frac{k_\alpha}{2}\int_0^t \sum_i\|\dot{q}_{i0}\|^2\,d\tau
		+ \frac{N\bar\sigma^2}{2k_\alpha} t. \label{eq5-3:V2_finite}
\end{equation}
Since $\dot{q}_{i0}$ is bounded on $[0,T^*]$ (from $s_i$ bounded and $q_{i0}$ continuous),
$V_2(t)$ is finite for every finite $t$.  In particular, $V_2(T^*)$ is finite.

Now suppose, for contradiction, that two agents $i$ and $j$ collide at some time $t_c \le T^*$,
i.e.~$\|q_{ij}(t_c)\| = d$.  Then the APF integral term diverges:
\begin{equation}
	\int_{\|q_{ij}(t_c)\|}^{\mu}\alpha(s)\,ds
	= \int_{d}^{\mu}\alpha(s)\,ds = \infty,
\end{equation}
which implies $V_2(t_c) = \infty$, contradicting the finiteness of $V_2$ on $[0,T^*]$.
Hence no collision occurs on $[0,T^*]$, and $\|q_{ij}(t)\| > d$ for all $t \in [0,T^*]$,
$i \neq j$.

\medskip\noindent\textbf{Step 5 (Post-sliding clean dynamics).}

For $t \ge T^*$, $s_i \equiv 0$ and $\dot{s}_i \equiv 0$.  From \eqref{eq5-3:qddot_transient},
\begin{equation}
	\ddot{q}_{i0} = -k_\alpha \dot{q}_{i0} - k_\alpha q_{i0} + \phi_i,
	\qquad \forall\, t \ge T^*. \label{eq5-3:clean_dyn}
\end{equation}
This is the same potential-driven dynamics as in Controller 2, but obtained without ever
computing $\dot{\phi}_i$.

\medskip\noindent\textbf{Step 6 (Velocity convergence --- P3).}

For $t \ge T^*$, differentiate $V_2$ along \eqref{eq5-3:clean_dyn}.  Substituting
\eqref{eq5-3:clean_dyn} into \eqref{eq5-3:V2_dot_transient} and noting $\dot{s}_i \equiv 0$,
\begin{equation}
	\dot{V}_2 = -k_\alpha\sum_{i\in\mathcal{N}} \|\dot{q}_{i0}\|^2 \le 0,
	\qquad \forall\, t \ge T^*. \label{eq5-3:V2_dot_clean}
\end{equation}
Hence $V_2(t)$ is nonincreasing and bounded below by zero, so it converges to a finite limit.
To apply Barbalat's lemma, we verify that $\ddot{V}_2$ is uniformly bounded.
Differentiating \eqref{eq5-3:V2_dot_clean},
\begin{equation}
	\ddot{V}_2 = -2k_\alpha\sum_{i\in\mathcal{N}} \dot{q}_{i0}^T \ddot{q}_{i0}.
\end{equation}
From \eqref{eq5-3:clean_dyn}, $\ddot{q}_{i0}$ is expressed in terms of $\dot{q}_{i0}$, $q_{i0}$
and $\phi_i$.  Since $\dot{q}_{i0}$ and $q_{i0}$ are bounded (from $V_2$ bounded) and $\phi_i$
is bounded (collision avoidance keeps $\|q_{ij}\| > d$, so $\alpha(\|q_{ij}\|)$ is bounded),
$\ddot{q}_{i0}$ is bounded, hence $\ddot{V}_2$ is bounded.  Barbalat's lemma then gives
$\dot{V}_2 \to 0$, i.e.~$\|\dot{q}_{i0}\| \to 0$ for every $i$.  Thus
$\lim_{t\to\infty} \|\dot{q}_i(t) - \dot{q}_0(t)\| = 0$, establishing P3.

\medskip\noindent\textbf{Step 7 (Convex-hull fencing --- P1).}

Define the centroid error $\tilde{q} = \frac1N\sum_{i\in\mathcal{N}} q_{i0}$.
Averaging \eqref{eq5-3:clean_dyn} over $i$ and using $\sum_i \phi_i = 0$,
\begin{equation}
	\ddot{\tilde{q}} = -k_\alpha \dot{\tilde{q}} - k_\alpha \tilde{q},
	\qquad \forall\, t \ge T^*. \label{eq5-3:centroid}
\end{equation}
This is a linear second-order system with characteristic polynomial
$\lambda^2 + k_\alpha \lambda + k_\alpha$.  For $k_\alpha > 0$, both roots have negative real
parts, so the origin $(\tilde{q},\dot{\tilde{q}}) = (0,0)$ is exponentially stable.
Consequently, $\tilde{q}(t) \to 0$ as $t \to \infty$.

Since the centroid $\frac1N\sum_{i=1}^N q_i(t) = q_0(t) + \tilde{q}(t)$ lies in the convex hull
$\operatorname{co}(q(t))$ by definition, we have
\begin{equation}
	\operatorname{dist}\bigl(q_0(t), \operatorname{co}(q(t))\bigr)
	\le \|\tilde{q}(t)\| \to 0,
\end{equation}
which establishes P1.

\medskip\noindent\textbf{Step 8 (Collision avoidance for all $t \ge 0$ --- P2).}

For $t \ge T^*$, \eqref{eq5-3:V2_dot_clean} gives $\dot{V}_2 \le 0$, so
$V_2(t) \le V_2(T^*) < \infty$.  The same barrier argument as in Step 4 applies:
if $\|q_{ij}(t)\| = d$ for some $t > T^*$, then $V_2(t) = \infty$, contradicting the
finiteness of $V_2(t)$.  Combined with Step 4, we conclude $\|q_{ij}(t)\| > d$ for all
$t \ge 0$ and all $i \neq j$, establishing P2.

\medskip\noindent\textbf{Remark on the PE condition.}

Throughout the proof, the PE condition is never required.  The finite-time convergence of
$s_i$ (Step 1) relies only on the robustness gain $\rho_i$ dominating the parametric
uncertainty $Y_i\tilde{\Theta}_i$, which is guaranteed by the projection operator and the
bounds on $\|\Theta_i\|$ and $\|Y_i\|$.  The post-sliding analysis (Steps 5--8) is purely
deterministic and independent of parameter convergence; indeed, $\tilde{\Theta}_i$ need not
tend to zero.  This is a key advantage over standard adaptive controllers that rely on PE
for asymptotic tracking.

\hfill$\square$

\begin{remark}
	Controller 5-3 replaces the gain condition $\rho_i \ge \|Y_i\tilde{\Theta}_i\| + \eta_i$
	(of Controller 5-1) with the two natural bounds in Assumption~(ii):
	$\|\Theta_i\| \le \Theta_{i,\max}$ and $\|Y_i\| \le k_Y$. The projection
	\eqref{eq5-3:proj} guarantees that $\|\tilde{\Theta}_i\| \le 2\Theta_{i,\max}$
	for all $t$, so the robust gain $\rho_i = k_Y \cdot 2\Theta_{i,\max} + \eta_i$ is a
	constant that can be computed once the physical parameter ranges and the operating
	envelope of the EL system are known.  The bound on $Y_i$ is not tied to the
	collision-avoidance proof; it follows from the continuity of $Y_i$ (Assumption~A3)
	and the boundedness of the EL system's physically realisable states (joint limits,
	bounded velocities, bounded target trajectory).  Hence there is no circularity with
	the collision-avoidance argument in Step~4.

	The price of the projection approach is the same chattering issue as in Controller
	5-1 (addressed by replacing $\mathrm{sgn}(s_i)$ with $\mathrm{sat}(s_i/\delta_i)$),
	and the need for the projection operator \eqref{eq5-3:proj} in the adaptation law.
\end{remark}